\MakeAbstract{
    本文设计并实现了一款轻量级实时数字协作白板系统 QuickBoard，面向课堂讨论与远程协作场景，强调无需注册、秒级创建/加入房间的使用体验。系统前端基于 Vue3 与 Fabric 实现可编辑画布，提供选择、画笔、形状、橡皮、文本与导出等基础能力；后端基于 Node.js 与 Socket.IO 实现房间管理与实时消息通信，并采用服务端权威裁决与增量同步策略减少冲突与回滚；同时接入独立 OCR 服务完成数学公式识别，并在前端以 KaTeX 渲染与可编辑弹层实现“识别—确认—插入—二次编辑”的闭环体验。本文围绕需求分析、系统设计与关键实现展开，给出核心模块的数据流与接口定义，并通过单元测试、构建验证与典型协作场景测试对系统可用性与稳定性进行评估。结果表明，该系统能够满足轻量协作白板的主要需求，并为后续扩展更丰富的图形组件与协作能力提供了可扩展的实现基础。
}{在线白板，实时协作，Socket.IO，OCR，KaTeX}

\MakeAbstractEng{
    This thesis designs and implements QuickBoard, a lightweight real-time digital collaborative whiteboard for classroom discussions and remote teamwork. It emphasizes instant room creation/joining without sign-up. The frontend is built with Vue 3 and Fabric to provide essential drawing capabilities such as selection, pen, shapes, eraser, text, and export. The backend is implemented with Node.js and Socket.IO for room management and real-time messaging, and adopts a server-authoritative decision with incremental synchronization to reduce conflicts and rollback issues. In addition, a standalone OCR service is integrated to recognize mathematical expressions, and KaTeX-based rendering with an editable workflow enables a complete loop of recognition, confirmation, insertion, and further editing. This thesis presents requirement analysis, system design, and key implementations with clarified data flows and interfaces, and evaluates availability and stability through unit tests, production builds, and scenario-based collaborative tests. The results show that the system meets the core needs of a lightweight collaborative whiteboard and provides a scalable foundation for future extensions.
}{online whiteboard, real-time collaboration, Socket.IO, OCR, KaTeX}
